\section{Results}
\label{sec:results}

We report strategy comparisons, the loop trajectory, multi-objective trade-offs, endpoint ablation, and cross-split robustness.

\subsection{Strategy Comparison}
\label{sec:strategy_comparison}

Table~\ref{tab:strategies} and Figure~\ref{fig:baseline_comparison} compare all eight strategies on the validation split ($k = 20$).
Bootstrap 95\% confidence intervals are computed over 1{,}000 resamples.

\begin{table}[t]
\centering
\caption{Strategy comparison on the validation split ($k = 20$, $n = 539$ candidates). All metrics are means across three oracle definitions. 95\% CIs from 1{,}000 bootstrap resamples. Best point estimate in \textbf{bold}.}
\label{tab:strategies}
\small
\begin{tabular}{@{}lcccc@{}}
\toprule
\textbf{Strategy} & \textbf{Top-$k$ Enrichment} & \textbf{NDCG} & \textbf{Hypervolume} & \textbf{Feasible Frac.}  \\
\midrule
Random & 0.037 \scriptsize{[0.000, 0.117]} & 0.536 \scriptsize{[0.444, 0.631]} & 0.364 & 0.126 \\
Activity only & 0.110 \scriptsize{[0.050, 0.200]} & 0.713 \scriptsize{[0.634, 0.798]} & 0.417 & 0.000 \\
Toxicity exclusion & 0.238 \scriptsize{[0.133, 0.367]} & 0.807 \scriptsize{[0.749, 0.863]} & 0.409 & \textbf{0.982} \\
Rule only & 0.396 \scriptsize{[0.250, 0.567]} & 0.865 \scriptsize{[0.804, 0.928]} & 0.431 & 0.450 \\
NSGA-II & 0.444 \scriptsize{[0.267, 0.617]} & 0.821 \scriptsize{[0.751, 0.896]} & \textbf{0.584} & 0.523 \\
Weighted sum & 0.609 \scriptsize{[0.483, 0.733]} & 0.976 \scriptsize{[0.955, 0.990]} & 0.554 & 0.679 \\
Random wt.\ search & 0.611 \scriptsize{[0.483, 0.734]} & 0.954 \scriptsize{[0.903, 0.982]} & 0.547 & 0.671 \\
\textbf{Agent improved} & \textbf{0.650} \scriptsize{[0.533, 0.767]} & \textbf{0.982} \scriptsize{[0.966, 0.992]} & 0.558 & 0.707 \\
\bottomrule
\end{tabular}
\end{table}

\begin{figure}[t]
\centering
\includegraphics[width=\linewidth]{figures/fig1_baseline_comparison.pdf}
\caption{Strategy comparison on the validation split ($k = 20$). Three grouped bars per strategy show top-$k$ enrichment, NDCG, and feasible fraction. Error bars are bootstrap 95\% CIs. The agent-improved policy leads on both top-$k$ enrichment and NDCG. Toxicity exclusion achieves the highest feasible fraction (0.982) but ranks poorly on enrichment (0.238), illustrating the multi-objective trade-off.}
\label{fig:baseline_comparison}
\end{figure}

The agent-improved policy achieves the highest point estimate on both top-$k$ enrichment (0.650) and NDCG (0.982), and the second-highest feasible fraction (0.707).
No other strategy leads on all three metrics simultaneously.

\textbf{Agent improved vs.\ NSGA-II.}
The strongest comparison is against NSGA-II, the standard multi-objective optimization baseline: 0.650 vs.\ 0.444 top-$k$ enrichment, with substantially separated 95\% CIs ([0.533, 0.767] vs.\ [0.267, 0.617]).
NSGA-II achieves the highest hypervolume (0.584), reflecting its design goal of spreading solutions across the Pareto front.
However, for candidate triage where the objective is to identify the top 20 candidates rather than a diverse frontier, scalarization-based policies are better suited.

\textbf{Agent improved vs.\ random weight search.}
The agent outperforms best-of-1{,}000 random weight search (0.650 vs.\ 0.611 top-$k$ enrichment), demonstrating that directed policy search adds value beyond blind exploration of the same policy class.
The agent discovered structural innovations (consensus voting, reciprocal rank fusion, learned reranking) that random weight sampling cannot access.

\textbf{Agent improved vs.\ weighted sum.}
The improvement over equal-weight scalarization (0.650 vs.\ 0.609) is real but the bootstrap CIs overlap: [0.533, 0.767] vs.\ [0.483, 0.733].
The improvement is not statistically significant at 95\% confidence on this single validation split.
However, the improvement is consistent across 10 independent splits ($0.650 \pm 0.034$ vs.\ $0.585 \pm 0.050$; Section~\ref{sec:cross_split}).

\textbf{Feasibility trade-off.}
Toxicity exclusion achieves the highest feasible fraction (0.982) but the lowest enrichment among non-trivial strategies (0.238), finding safe peptides while missing the best candidates.
Activity only achieves 0.000 feasible fraction.
The agent-improved policy balances all three metrics.

\subsection{Loop Trajectory}
\label{sec:trajectory}

Figure~\ref{fig:loop_trajectory} shows the autoresearch progress over 100 experiments on the loop's held-out validation split ($n = 533$; the bootstrap and multi-split evaluations in Sections~\ref{sec:strategy_comparison}--\ref{sec:cross_split} use the final split with $n = 539$).
Of 100 experiments, 12 are kept (11 improvements plus the initial baseline), yielding a keep rate of 11\%.
The harmonic mean of top-$k$ enrichment and NDCG on the loop's split progresses from 0.802 (experiment 1) to 0.828 (experiment 94).

\begin{figure}[t]
\centering
\includegraphics[width=\linewidth]{figures/fig2_loop_trajectory.pdf}
\caption{Autoresearch loop trajectory over 100 experiments. Each point is one experiment; green triangles indicate kept improvements, gray circles indicate discarded regressions. The staircase line tracks the running best. The agent progresses through three structural phases: weight tuning (experiments 1--10), two-stage reranking (11--31), and oracle consensus voting with learned reranking (32--94, with no further improvements through experiment 100).}
\label{fig:loop_trajectory}
\end{figure}

The 12 kept experiments reveal a structural progression across three phases:

\textbf{Phase 1: Weight tuning} (experiments 1--10).
The agent adjusts scalarization weights, increasing activity weight from 0.45 to 0.50 and toxicity weight from 0.40 to 0.50.
Top-$k$ enrichment improves from 0.683 to 0.700.

\textbf{Phase 2: Two-stage reranking} (experiments 11--31).
The agent discovers that selecting a top-$k$ shortlist via weighted sum and reranking by a nonlinear threshold-gated score outperforms direct scalarization.
It then introduces reciprocal rank fusion \citep{cormack2009rrf} across multiple heuristic scores (threshold-gated, rank product, Pareto).

\textbf{Phase 3: Oracle consensus voting with learned models} (experiments 32--94).
The agent introduces a consensus voting mechanism: candidates appearing in at least two of the three heuristic top-20 shortlists (threshold-gated, rank product, Pareto) are prioritized.
This is equivalent to majority-vote ensemble selection across three independent heuristic scores.
Top-$k$ enrichment reaches 0.717.
Subsequent experiments integrate LambdaMART \citep{ke2017lightgbm,burges2010ranknet} and MLP rerankers trained on oracle-aligned labels, with the final policy using a per-oracle MLP bottleneck score for shortlist reranking.

Table~\ref{tab:keeps} in Appendix~\ref{sec:supp_tables} lists all 12 kept experiments with their metrics and descriptions.

\subsection{Multi-Objective Trade-Offs}
\label{sec:multiobjective}

Figure~\ref{fig:pareto_front} visualizes the multi-objective trade-off by plotting the top-20 candidates selected by three representative strategies in the (activity, $1 - \text{toxicity}$) plane across the three oracle definitions.

\begin{figure}[t]
\centering
\includegraphics[width=\linewidth]{figures/fig3_pareto_front.pdf}
\caption{Pareto front visualization. Each panel shows one oracle's top-20 selected by three strategies in the (activity, $1 - \text{toxicity}$) plane. Activity only (left) clusters in the high-activity/low-safety corner. The agent-improved policy (right) covers a broader region of the Pareto front while maintaining competitive activity.}
\label{fig:pareto_front}
\end{figure}

Activity-only selection concentrates in the high-activity region regardless of toxicity, yielding 0.000 feasible fraction.
The agent-improved policy distributes selections more broadly across the (activity, safety) plane, achieving 0.707 feasible fraction while maintaining the highest enrichment.

\subsection{Endpoint Ablation}
\label{sec:ablation}

To quantify the contribution of each endpoint, we ablate by setting each endpoint's values to their pool mean and re-evaluating (Table~\ref{tab:ablation}, Figure~\ref{fig:ablation_heatmap}).

\begin{table}[t]
\centering
\caption{Endpoint ablation: change in top-$k$ enrichment ($\Delta$TopK) and NDCG ($\Delta$NDCG) when each endpoint is neutralized (set to pool mean). Positive $\Delta$TopK indicates the strategy \emph{improves} when an endpoint is removed (i.e., the endpoint was hurting). Evaluated on the validation split.}
\label{tab:ablation}
\small
\begin{tabular}{@{}llrr@{}}
\toprule
\textbf{Strategy} & \textbf{Ablated endpoint} & \textbf{$\Delta$TopK} & \textbf{$\Delta$NDCG} \\
\midrule
\multirow{4}{*}{Weighted sum}
  & Activity & $+$0.133 & $+$0.013 \\
  & Toxicity & $\phantom{+}$0.000 & $+$0.003 \\
  & Stability & $+$0.033 & $-$0.014 \\
  & Developability & $+$0.033 & $-$0.010 \\
\midrule
\multirow{4}{*}{Agent improved}
  & Activity & $+$0.167 & $+$0.003 \\
  & Toxicity & $+$0.117 & $-$0.002 \\
  & Stability & $+$0.083 & $-$0.006 \\
  & Developability & $+$0.033 & $-$0.006 \\
\midrule
\multirow{4}{*}{Rule only}
  & Activity & $+$0.183 & $+$0.065 \\
  & Toxicity & $+$0.067 & $-$0.025 \\
  & Stability & $-$0.200 & $-$0.029 \\
  & Developability & $+$0.067 & $+$0.008 \\
\bottomrule
\end{tabular}
\end{table}

\begin{figure}[t]
\centering
\includegraphics[width=0.85\linewidth]{figures/fig4_ablation_heatmap.pdf}
\caption{Endpoint ablation heatmap. Each cell shows the change in top-$k$ enrichment when that endpoint is set to its pool mean. Dropping stability causes the largest degradation for rule only ($-0.200$), while the agent-improved policy is more robust to individual endpoint ablation.}
\label{fig:ablation_heatmap}
\end{figure}

The ablation reveals that the effect of removing an endpoint is strategy-dependent.
Removing activity improves all three strategies' enrichment on this split (the oracles partially decouple from raw activity), while removing stability degrades rule only severely ($\Delta$TopK $= -0.200$) because rule only sorts primarily by stability.
The agent-improved policy shows moderate sensitivity across all four endpoints without catastrophic dependence on any single one.

\subsection{Cross-Split Robustness and Statistical Significance}
\label{sec:cross_split}

To assess robustness beyond a single validation split, we evaluate all strategies across 10 independent random splits (Figure~\ref{fig:multi_split}; full table in Appendix~\ref{sec:supp_multi_split}).

\begin{figure}[t]
\centering
\includegraphics[width=0.85\linewidth]{figures/fig6_multi_split_boxplot.pdf}
\caption{Cross-split robustness: top-$k$ enrichment across 10 random train/val/test splits. Agent improved ($0.650 \pm 0.034$) shows a higher mean and tighter variance than weighted sum ($0.585 \pm 0.050$), with non-overlapping interquartile ranges. See Figure~\ref{fig:cross_split_paired} for the paired per-split comparison.}
\label{fig:multi_split}
\end{figure}

The agent-improved policy achieves the highest mean top-$k$ enrichment ($0.650 \pm 0.034$) across all 10 splits, with the lowest standard deviation among the top three strategies.
Across 10 independent splits, agent improved beats weighted sum on 8 of 10 (mean gain $+0.065$; the two losses are small, $\Delta \leq 0.033$, while the eight wins are substantially larger, mean $\Delta = +0.088$).
A Wilcoxon signed-rank test on the paired per-split deltas gives $p = 0.004$ (one-sided), which is the appropriate test because it uses the magnitude of the wins, not only their sign.
Figure~\ref{fig:cross_split_paired} visualizes the paired comparison: eight of ten splits slope upward from weighted sum to agent improved (green), and the two downward slopes (red) are small in magnitude.

\begin{figure}[t]
\centering
\includegraphics[width=\linewidth]{figures/fig3c_cross_split_paired.pdf}
\caption{Paired per-split comparison, weighted sum vs.\ agent improved, on the benchmark's primary metric (mean top-$k$ enrichment across three oracles; $k = 20$). Left: slope graph with one line per split seed (100--109); green lines indicate agent improved wins ($8/10$), red lines indicate losses ($2/10$). Right: per-split delta $\Delta = \text{agent} - \text{weighted sum}$; the two losses (seeds 100, 102) are small ($\Delta \geq -0.033$), while the eight wins spread to $+0.15$. Wilcoxon signed-rank $p = 0.004$ (one-sided).}
\label{fig:cross_split_paired}
\end{figure}

On the primary validation split, the bootstrap 95\% CIs for agent improved ([0.533, 0.767]) and weighted sum ([0.483, 0.733]) overlap; the improvement is not statistically significant at 95\% confidence on a single split.
The CIs are inherently wide because top-$k$ enrichment at $k = 20$ is a discrete metric with limited resolution.
Against NSGA-II, the single-split CIs are substantially separated ([0.533, 0.767] vs.\ [0.267, 0.617]; point estimate gap of 0.206).

Overall, the agent-improved policy outperforms all baselines on point estimates, with large margins against NSGA-II and weaker baselines (substantially separated CIs).
The margin over weighted sum is practically meaningful and cross-split-consistent, but bootstrap CIs overlap on any single split.
